\documentclass[11pt,a4paper]{article}
\usepackage[utf8]{inputenc}
\usepackage[T1]{fontenc}
\usepackage[ngerman]{babel}
\usepackage[margin=2.4cm]{geometry}
\usepackage{booktabs}
\usepackage{xcolor}
\usepackage[colorlinks=true,urlcolor=blue!60!black,linkcolor=black]{hyperref}
\usepackage{parskip}
\usepackage{enumitem}
\setlist{nosep,leftmargin=*}
\emergencystretch=3em

\newcommand{\code}[1]{\texttt{\small #1}}

\title{AWI-ESM3-4-2-veg-HR, CMIP7 piControl\\[2mm]
\large Stand nach \code{cli119}: \code{latitude}/\code{longitude}, und \code{\#81} auf dem ganzen Jahr}
\author{Jan Streffing, AWI}
\date{13. September 2026}

\begin{document}
% babel-ngerman macht " aktiv, was in Code-Beispielen zuschlaegt. Die Shorthands
% werden erst bei \begin{document} scharfgeschaltet, also muss das hierhin.
\shorthandoff{"}
\maketitle

\section{Kurzfassung}

Danke fuer den Lauf mit den Koordinatenchecks auf \code{cli118}. Der eine Befund,
die Namen der Hilfskoordinaten, ist behoben. \code{cli119} ist das Jahr 1851 mit
529 Dateien, geschrieben mit dieser Korrektur. Wir haben es komplett mit deinem
\code{\#81} in \code{5a4239e} geprueft, also dem Stand nach deinem Lauf:

\begin{center}
\begin{tabular}{lrr}
\toprule
\textbf{HIGH} & \textbf{\code{cli118}, dein Lauf} & \textbf{\code{cli119} mit \code{\#81}} \\
\midrule
\code{wcrp\_cmip7} \code{COORD011}, Namen der Hilfskoordinaten & 455 Dateien & 0 \\
\code{wcrp\_cmip7} \code{TIME001}                  & 0 & 0 \\
\code{aicc} \code{vegtype} auf \code{landCoverFrac} & 1 & 1 \\
cf Appendix A, \code{formula\_terms} auf \code{lev\_bnds} & 9 & 9 \\
\bottomrule
\end{tabular}
\end{center}

\code{wcrp\_cmip7} meldet damit auf dem ganzen Jahr kein HIGH mehr. Was bleibt,
liegt in der Tabelle und im cf-Checker, nicht in den Dateien.

Dazu ein Fehler bei uns, den kein Checker sehen konnte: die dekadischen Dateien
enthielten bis einschliesslich \code{cli119} nur ein Jahr. Das ist behoben,
Abschnitt~\ref{sec:dec}, mit zwei Fragen an dich.

\section{\code{latitude} und \code{longitude}}

Zwei Tabellen benennen die horizontalen Koordinaten. \code{CMIP7\_coordinate.json}
gilt fuer die Achsen regulaerer Gitter, \code{out\_name} \code{lat}/\code{lon}.
Darauf hatten wir in der dritten Runde 41 Dateien von \code{latitude(latitude)} auf
\code{lat(lat)} umgestellt. \code{CMIP7\_grids.json} gilt fuer die Gittervariablen
kurvilinearer und unstrukturierter Gitter, \code{out\_name}
\code{latitude}/\code{longitude}. Dort hatten unsere Hilfskoordinaten einfach den
Namen behalten, den XIOS oder das Mesh mitbrachte. Keiner der bisherigen Checker hat
darauf geschaut.

Hilfskoordinaten heissen jetzt \code{latitude}/\code{longitude}, und jedes
\code{coordinates}-Attribut zieht mit. Regulaere Gitter behalten \code{lat(lat)}
und \code{lon(lon)}, die Umstellung aus Runde drei gilt also weiter.

\begin{center}
\begin{tabular}{lrl}
\toprule
\textbf{Gitter} & \textbf{Dateien} & \textbf{in \code{cli119}} \\
\midrule
\code{g122} & 321 & \code{latitude(ncells)} bzw.\ \code{latitude(cell)} \\
\code{g130} & 125 & \code{latitude(ncells)} \\
\code{g132} &   9 & \code{latitude(elem)} \\
\code{g129}, \code{g239} & 45 & \code{lat(lat)}, unveraendert \\
\code{g010} & 29 & ohne horizontale Koordinaten \\
\bottomrule
\end{tabular}
\end{center}

\code{g132} stand nicht in deiner Liste, war aber genauso betroffen, die neun
Dateien sind mit umgestellt.

\section{\code{TIME001}}

In deinem Lauf auf \code{cli118} war \code{TIME001} schon verschwunden. Seitdem hat
\code{\#81} die Pruefung noch einmal umgebaut, die Intervallgrenzen werden jetzt
aus dem Dateinamen neu erzeugt statt den deklarierten Grenzen zu vertrauen. Genau
unsere 42 subdaily- und dekadischen Dateien haengen daran, deshalb wollten wir das
nachmessen. Auch mit \code{5a4239e} bleiben alle 42 ohne Befund.

Auf \code{cc-plugin-wcrp\#80} hat \code{Ayoubnac1} geschrieben, dass mehrere Dekaden in einer
Datei unter \code{proleptic\_gregorian} noch um zwoelf Stunden daneben liegen. Das
betrifft uns nicht, jede unserer \code{dec}-Dateien umfasst genau eine Dekade.
Was in dieser Dekade stand, war allerdings falsch, dazu der naechste Abschnitt.

\section{Dekadenmittel: ein Fehler bei uns}
\label{sec:dec}

Wir cmorisieren jahresweise, ein Modelljahr pro Job, und jedes Eingabemuster wird
auf dieses Jahr gepinnt. Das galt auch fuer die neun \code{dec}-Variablen
(\code{thetao}, \code{so}, \code{tauuo}, \code{tauvo}, \code{thkcello}, \code{volo},
\code{masscello}, \code{volcello}, \code{masso}). Die Grenzen haben wir korrekt als
1850 bis 1860 geschrieben, gemittelt wurde aber nur das jeweils verarbeitete Jahr.
Jedes weitere Jahr derselben Dekade haette dieselbe Datei mit einem anderen
Einjahresmittel ueberschrieben. Strukturell sind diese Dateien in Ordnung, deshalb
ist es in keiner QC aufgefallen. Die neun \code{dec}-Dateien aus \code{cli117} bis
\code{cli119} sind nicht zu verwenden.

\textbf{Jetzt.} Jeder Jahresjob traegt die \code{dec}-Regeln weiter mit. Ein erster
Schritt zaehlt das Startjahr des Laufs jedes Mal neu aus den Eingabedateien und
zaehlt Dekaden ab diesem Jahr. Die Regel schreibt nur in dem Jahr, das eine Dekade
schliesst, und liest dann alle zehn Jahre; in jedem anderen Jahr endet sie ohne
Ausgabe. Fehlt in einer schliessenden Dekade ein Jahr, ist das ein Fehler, kein
Uebergehen. Eine unvollstaendige letzte Dekade wird nicht geschrieben.

\textbf{Nachgemessen} fuer 1858 und 1859:

\begin{itemize}
\item 1858: alle neun Regeln enden ohne Ausgabe.
\item 1859: neun Dateien, \code{time} 1855-01-01, Grenzen 1850-01-01 bis
  1860-01-01, jeweils zehn Eingabedateien gelesen.
\item \code{tauuo} an 2000 zufaelligen Zellen gegen das Mittel der 120 Monatswerte:
  Abweichung $10^{-5}$, also \code{float32}-Genauigkeit. Gegen 1859 allein: 115\,\%.
\item \code{\#81} in \code{5a4239e}: kein HIGH auf den neun Dateien.
\end{itemize}

\textbf{Warum ab Laufbeginn und nicht nach Kalenderdekaden.} Weder die CV
(\code{dec}: \glqq Decadal samples\grqq, Intervall zehn Jahre) noch die CMIP7-Guidance
legen fest, wo eine Dekade beginnt. \code{\#81} ist es ebenfalls egal, solange die
Zelle zehn Jahre lang ist: 1850--1860, 1851--1861 und 2024--2034 bestehen,
2024--2030 und 2040--2046 werden als HIGH gemeldet (\glqq does not match regular
10-year cells\grqq). In \code{/pool/data/CMIP6} ist es uneinheitlich: AWI-CM-1-1-MR
und AWI-ESM-1-1-LR zaehlen 1851--1860 und lassen angebrochene Dekaden weg,
IPSL-CM6A-LR nimmt fuer \code{historical} 1850--1859, fuer \code{ssp585} aber
Dekaden ab 2015, MIROC6 hat eine Zelle 2010--2014 veroeffentlicht.

\textbf{Zwei Fragen dazu:}
\begin{enumerate}
\item Ist es fuer euch in Ordnung, dass eine angebrochene letzte Dekade fehlt, im
  \code{piControl} also 2040--2045, statt als kuerzere Zelle zu erscheinen?
\item Sollen Szenarien ihre Dekaden ab Experimentbeginn zaehlen, bei einem Start
  2024 also 2024--2033, oder nach Kalenderdekaden ab 2030? Ab Beginn geht nichts
  verloren, nach Kalender liegen \code{historical} und Szenario auf demselben Raster.
\end{enumerate}

\section{Was uebrig bleibt}

\begin{center}
\begin{tabular}{rll}
\toprule
\textbf{Anzahl} & \textbf{Befund} & \textbf{liegt bei} \\
\midrule
1 HIGH     & \code{vegtype}, Singular \code{needleleaf\_evergreen\_tree} & \code{CMIP7\_DReq\_Content\#29} \\
9 HIGH     & cf Appendix A gegen CF 1.7 \S7.1, \code{formula\_terms} & \code{compliance-checker} \\
492 MEDIUM & \glqq Registry has no value for key 'cell\_measures'\grqq{} & \code{WCRP-universe\#334} \\
\bottomrule
\end{tabular}
\end{center}

\textbf{Eine Korrektur zur fuenften Antwort.} Dort stand, dass die 492 MEDIUM mit
\code{\#81} vermutlich verschwinden, weil der PR die Kompatibilitaet von
\code{units} und \code{cf\_units} behandelt. Das stimmt nicht: mit \code{5a4239e}
sind es weiterhin 492. Die Universe-Datenbank 3.2.0 fuehrt fuer
\code{known\_branded\_variable} schlicht kein \code{cell\_measures}, und das kann der
Checker nicht ausgleichen. Laut \code{WCRP-universe\#334} kommt die Korrektur mit
dem QA/QC-Release fuer die Koordinatenchecks.

Ein MEDIUM mehr als ohne \code{\#81} gibt es auf den beiden \code{time4}-Dateien:
\code{\#81} meldet den Tippfehler \code{Satistics} aus der Tabelle, der schon in
deinem \code{CMIP7\_DReq\_Content\#31} steht.

\section{Zu deinen weiteren Punkten}

Zum Gitterregister: solange \code{grid\_type} nicht aus EMD in die CV kommt,
behalten wir unser eigenes Register fuer \code{aicc}. Wenn \code{\#81} ueber
\code{grid\_label} auf die Topologie kommt, pruefen wir, ob unsere Labels dort
hinterlegt sind. \code{g239} ist inzwischen auch auf \code{main} von
\code{CMIP7-CVs} angekommen.

Die Konsistenzbefunde zwischen Datensaetzen aus \code{esgf-qa} lassen wir liegen,
bis du die Gruppen neu gebildet hast.

\section{Wie gemessen wurde}

Die QC im Lauf selbst ist gegenueber \code{cli118} unveraendert: cf 0 HIGH, 128
MEDIUM, \code{wcrp\_cmip7} \code{develop} \code{ba2de16} 42 HIGH und 613 MEDIUM,
\code{aicc} 1 HIGH. Die Umbenennung kostet also mit den heute veroeffentlichten
Checkern nichts.

Da keine der QC-Suiten die Werte selbst prueft, laeuft bei uns zusaetzlich ein
Abgleich jeder Datei mit Literaturbereichen fuer Minimum, Mittel und Maximum. Auf
\code{cli119} bestehen 307 von 529 Dateien, 169 liegen knapp ausserhalb, 53 deutlich.
49 der 53 fielen schon im August auf, die Ursachen gehen wir gerade durch.

Die Pruefung mit \code{\#81} lief danach separat ueber alle 529 Dateien, mit dem
PR-Stand \code{5a4239e} in einem eigenen Checkout, \code{compliance-checker}
6.1.0 mit aktivierten Appendix-A-Checks, \code{cc-plugin-aicc} 0.2.0 mit unserem
Register, esgvoc 6.1.0 mit \code{cmip7@2.3.0} und \code{universe@3.2.0}. Die
Dateien liegen unter
\code{/scratch/a/a270092/pycmor\_hr/cli119/cmorized/MIP-DRS7}, die Reports aus der
Pruefung mit \code{\#81} unter \code{/scratch/a/a270092/pycmor\_hr/cli119\_pr81}.

\end{document}
